\documentclass[12pt]{exam}

\usepackage{amssymb,amsmath,amsfonts,amsthm,graphicx}
\usepackage[top=0.3in, bottom=0.75in, left=0.5in, right=0.5in]{geometry}
\usepackage{tikz}
\DeclareGraphicsExtensions{.pdf}

\newtheoremstyle{problem}{\topsep}{\topsep}{}{}{\bfseries}{}{ }%
{\thmnumber{#2}\thmnote{ \bfseries (#3)}}
\theoremstyle{problem}
\newtheorem{p}{}

\firstpageheader%
{CS \liningnums{310}: Algorithms --- Graphs\\
Mudassir Shabbir, February 25, 2026}
{}
{Full name1:\enspace\makebox[2in]{\hrulefill}\\
Full name2:\enspace\makebox[2in]{\hrulefill}\\}

\runningheader{}{}{}

\frenchspacing
\unframedsolutions
\SolutionEmphasis{\sffamily}
\renewcommand{\solutiontitle}{}
\boxedpoints
\pointsinrightmargin
\marginbonuspointname{\textsc{up}}

\makeatletter
\def\clap#1{\hbox to 0pt{\hss#1\hss}}
\def\setup@point@toks{%
	\point@toks={%
		\rlap{\hskip-\@totalleftmargin
			\hskip\textwidth
			\hskip\@rightmargin
			\hskip-\rightpointsmargin
			\clap{\padded@point@block}%
		}%
		\global \point@toks={}%
	}%
}
\setlength{\rightpointsmargin}{.5in}
\makeatother

\extraheadheight{.25in}
\extrafootheight{-.5in}
\setlength{\marginparwidth}{1.5in}

\usepackage{xcolor}
\usepackage{pagecolor}
\usepackage[document]{ragged2e}

\begin{document}
	
\pagestyle{empty}
\noindent
\begin{tabular*}{\textwidth}{@{\extracolsep{\fill}} l r}
\textbf{CS 310: Algorithms -- Graphs} 
& \emph{Mudassir Shabbir, February 25, 2026}
\end{tabular*}

\vspace{0.4cm}

\noindent
\begin{tabular*}{\textwidth}{@{\extracolsep{\fill}} l l}
\textbf{Student ID 1 \& 2:}{\em (Do this in pairs)} \hrulefill 
&
\end{tabular*}

\vspace{0.3cm}
\hrule

\thispagestyle{myheadings}
\pagenumbering{gobble}


\vspace{0.5cm}
\smallskip
\textbf{A.} Recall that in Union by Rank, the root with smaller rank is attached under the root with larger rank. The rank only increases when two equal-rank trees are merged.
\textbf{Prove that Union by rank results in $O(\log n)$ cost per Find in Union-Find Data Structure.}

\vspace{0.5cm}

\textit{Hint:} Show that a tree of rank $r$ must contain at least $2^r$ nodes. Use this to bound the maximum possible height.

\vspace{7.5 cm}

\textbf{B.} \textbf{Design a decimal counter and perform amortized analysis of number of digits that need to be changed per increment (in class we did this for binary counter.)} The counter is in base 10, with digits $0$–$9$. When a digit becomes $9$, it resets to $0$ and generates a carry.

\vspace{0.5cm}

\textit{Hint:} Count how many times each digit position changes over $n$ increments. 
First compute the total number of digit changes across all $n$ increments, and then divide by $n$ to obtain the amortized cost per increment.



\end{document}